% 本文件由 LaTeX 排版 Agent 根据有效 translation.json 自动生成。
% author=Miscellaneous; work_id=1004; title=巴拉盎与约撒法特传奇

\chapter{巴拉盎与约撒法特传奇}

% structure: p0001 source=h1 target=omitted reason=page-title-already-rendered-by-parent

\Emph{阿里斯提德护教书，据巴拉盎与约撒法特史中所载。}\par

1. 我，哦，王，因天主的眷顾来到世界；当我观察天地、日月及其余万物时，我惊叹于它们井然的秩序。\par

当我看见宇宙及其中的一切皆因必然性而运动时，我认识到那推动和掌管者就是天主。\par

因为凡引起运动的，比那被运动的更强；凡掌管的，比那被掌管的更强。\par

那么，这位最初建立并至今掌管宇宙的同一存在——我肯定他就是天主，无始无终，不死而自足，超越一切情感与软弱，超越愤怒、遗忘、无知及其他。\par

万物也靠祂而存立。祂不需要祭献、奠酒或任何可见之物；但众人都需要祂。\par

2. 关于天主，我已尽我所能述说；接下来让我们论及人类，看看他们中谁参与真理，谁陷于错误。\par

因为在我们看来，哦，王，世上有三类人：即你们所敬拜的神灵的崇拜者、犹太人和基督徒。此外，那些敬拜多神的又分为三类，即加色丁人、希腊人和埃及人；因为这些人在对众多名号的尊神的服侍与敬拜中，成了其余民族的向导和导师。\par

3. 让我们看看他们中谁参与真理，谁陷于错误。\par

加色丁人因不认识天主，便迷失于元素，开始敬拜受造物胜过造物主。\par

他们为这些元素塑造了某些形状，称之为天、地、海、太阳、月亮及其他原始物体或光体的象征。他们将它们关闭在神龛中，敬拜它们，称它们为神，尽管他们不得不严密守卫它们，以防被贼偷走。他们没有认识到，守卫者比被守卫者更大，制造者比被制造者更大。因为若他们的神无力照管自身安全，又怎能保护他人？加色丁人敬拜无生命、无用的偶像，陷入了极大的错误。\par

我觉得奇怪的是，哦，王，他们所谓的哲学家竟全然没有观察到元素本身是会朽坏的。若元素是会朽坏且受制于必然性的，它们怎会是神？若元素不是神，那么为尊崇它们而造的像又怎会成为神？\par

4. 那么，哦，王，让我们继续论及元素本身，以表明它们不是神，而是可朽坏、可变的，是受造于虚无、依真天主——那不可毁坏、不变、不可见、且遍观万物、按己意改变万物者——的命令而产生的。关于元素，我还有什么可说的呢？\par

那些相信天是神的人错了。因为我们看见它因必然性而旋转运动，由许多部分构成，因而被称为有秩序的宇宙（\OriginalTerm{宇宙}{Kosmos}）。宇宙是一位设计者的构造；被构造之物有始有终。天及其光体因必然性而运动。群星以固定的间隔从一宫到另一宫列队运行，有些陨落，有些升起，按季节完成它们的行程，以划分夏冬，这是天主为它们所定的；它们服从本性不可逃避的必然性，不逾越各自的界限，与天上的秩序保持一致。由此可见，天不是神，而是天主的作为。\par

那些相信地是女神的人也错了。因为我们看见它受人滥用和蹂躏，被犁耕、被揉捏，变得毫无价值。若被火焚烧，它就失去生命；因为灰烬中什么也不会生长。此外，若雨水过多，它及其果实都会溶解。它还被人类及其他生物践踏；被谋杀者的血染红；被挖掘，填满尸体，成为死者的坟墓。面对这一切，不可能承认地是女神，而只是天主为人类使用所造的工程。\par

5. 那些相信水是神的人也错了。因为水也是为人类使用而造的，受人控制；它被玷污、毁坏，煮沸和染色时会发生变化；它被霜冻结，被血污染，并被用来洗涤一切不洁之物。因此，水不可能是神，而是天主的工程。\par

那些相信火是神的人也错了。因为火是为人类使用而造的，受人控制，可被从一个地方带到另一个地方，用于煮烤各种肉类，甚至焚烧尸体。它还能以多种方式被熄灭，借人力扑灭。所以不能承认火是神，而是天主的工程。\par

那些认为风的吹动是女神的人也错了。因为显然它受制于另一者；为人的缘故，天主设计了风，用于船舶运输、谷物运送及人的其他需要。它遵照天主的命令而起落，由此可见风的吹动不是女神，只是天主的工程。\par

6. 那些相信太阳是神的人也错了。因为我们看见它因必然性而运动、旋转，从一宫到另一宫，升起落下，为人类的使用而温暖植物和嫩芽。\par

此外，它与其他星辰共享一部分，并远小于天空；它的光会亏缺，且不服从自身的规律。因此可以断定，太阳不是神，而只是天主的化工。那些相信月亮是女神的人同样错误。因为我们看见它必然地运行、旋转，从一个宫位移到另一个宫位，为人类的益处而落下和升起；它小于太阳，有圆缺，也有蚀。因此可以断定，月亮不是女神，而只是天主的化工。\par

7. 那些相信人是神的人也错了。因为我们看见他必然地被推动，被迫成长，不愿老去却变老。他时而欢乐，时而因缺乏食物、饮水和衣服而悲伤。我们看见他易怒、嫉妒、贪欲、改变主意，并有许多软弱。他也被各种元素和动物以及永无休止的死亡以多种方式毁灭。因此不能承认人是神，而只是天主的化工。\par

因此，加色丁人随从自己的欲望而迷路，陷入了极大的错误。\par

因为他们崇拜可朽坏的元素和没有生命的偶像，却不知道是他们自己将这些造成为神。\par

8. 让我们转而看希腊人，看他们对天主是否有什么认识。希腊人虽然自称为智慧人，却证明比加色丁人更为迷惑，因为他们声称有许多神存在，有些是男神，有些是女神，在每一种情欲和每一种愚昧上都是老练的能手。[而希腊人自己又把他们描绘为奸夫、凶手、易怒、嫉妒、暴躁、弑父杀兄、小偷、强盗、跛脚、瘸腿、行巫术者、疯狂者。其中有些（据他们的说法）死了，有些被雷劈，成了人的奴隶，成为逃亡者，他们哀哭悲叹，为了邪恶可耻的目的把自己变成动物。]\par

因此，大王，希腊人所引进的这些故事是可笑、荒谬、不敬的，他们将不是神的东西冠以神之名，以迎合自己不圣洁的欲望，为的是有了这些恶习的保护者，他们可以行奸淫、抢劫、杀人和其他骇人听闻的事。因为如果他们的神做了这些事，他们为什么不该做呢？\par

因此，从这些误入歧途的做法，人类遭遇了频繁的战争、屠杀和痛苦的被掳。\par

9. 但是，如果我们有心逐一讨论他们的神，你们会看到那是多么荒谬；例如，他们提出\Person{克罗诺斯}是至高无上的神，且向他献祭自己的儿女。他由\Person{瑞亚}生了众多子女，并在疯狂中吞吃自己的后代。他们还说他被\Person{宙斯}阉割并将器官抛入海中，传说\Person{阿佛洛狄忒}由此而生。然后\Person{宙斯}捆绑了自己的父亲，把他投入\Place{塔耳塔罗斯}。你们看他们加于其神的错误和野蛮？那么，神难道可以被镣铐、被阉割吗？多么荒谬！哪个有理智的人会这样说呢？\par

接着又引进了\Person{宙斯}，他们说他是诸神之王，他把自己变成动物，以便奸淫凡间的女子。\par

因为他们声称他为了\Person{欧罗巴}变成一头公牛，为了\Person{达娜厄}变成金雨，为了\Person{勒达}变成一只天鹅，为了\Person{安提俄珀}变成一位萨提尔，为了\Person{塞墨勒}变成一道霹雳。然后由此生了许多儿女：\Person{狄俄倪索斯}、\Person{仄托斯}、\Person{安菲翁}、\Person{赫拉克勒斯}、\Person{阿波罗}、\Person{阿耳忒弥斯}、\Person{珀耳修斯}、\Person{卡斯托耳}、\Person{海伦}、\Person{波鲁克斯}、\Person{弥诺斯}、\Person{剌达曼堤斯}、\Person{萨耳珀冬}，以及被称为\Person{缪斯}的九个女儿。然后他们还提出关于\Person{伽倪墨得斯}的故事。\par

因此，大王，人类就模仿这一切，成为奸淫的男人和放荡的女人，并做了其他可怕的恶事，因为他们模仿他们的神。那么，一个神怎可能是奸夫、淫秽者或弑亲者呢？\par

10. 他们还随着他提出\Person{赫淮斯托斯}为神，说他是瘸子，挥舞锤子和钳子，以铁匠为生。\par

那么他贫困吗？但不可承认神是瘸子，并且还要依赖人类。\par

然后他们提出\Person{赫耳墨斯}为神，将他描绘为好色、偷窃、贪婪、行巫术（并且残废）、译解语言者。但不可承认这样一位是神。\par

他们还提出\Person{阿斯克勒庇俄斯}为神，说他是一位医生，调制药物、配制膏药以维持生计。因为他很贫穷。后来他们说他因\Person{拉刻岱蒙}之子\Person{廷达瑞俄斯}的事被\Person{宙斯}用霹雳击中，于是被杀。如果\Person{阿斯克勒庇俄斯}身为神，被雷击时连自己都救不了，又如何能拯救别人呢？\par

再者，\Person{阿瑞斯}被描绘为神，好斗、嫉妒、喜爱野兽和其他诸如此类的事物。最后在勾引\Person{阿佛洛狄忒}时，被年轻的\Person{厄洛斯}和\Person{赫淮斯托斯}捆绑。那么，一个受欲望支配、好战、被囚、行奸淫的人，怎会是神呢？\par

他们还声称\Person{狄俄倪索斯}是神，他举行夜宴，教人酗酒，拐带邻人之妻，发疯逃跑，最后被\Person{泰坦}杀死。如果\Person{狄俄倪索斯}在被杀时不能救自己，而且常发疯、醉于酒、逃亡，他怎会是神呢？\par

他们还声称\Person{赫拉克勒斯}酗酒、发疯，割断自己儿女的喉咙，然后被火烧死。那么，一个酗酒、杀子、被烧死的人怎会是神呢？或者，他连自己都救不了，又如何帮助别人呢？\par

11. 他们也把\Person{阿波罗}描绘为忌妒的神，此外还是弓箭的主人，有时是竖琴和笛子的主人，并为报酬给人类占卜？那么他非常贫困吗？但不可承认神会有缺乏、忌妒，并且是个弹唱乐师。\par

他们又将\Person{阿耳忒弥斯}描述为他的姐妹，她是一名女猎人，带着弓箭和箭袋；她独自在山丘上带着狗漫游，猎杀公鹿或野猪。那么这样一个带着狗打猎漫游的女人，又怎么会是神呢？\par

他们甚至断言\Person{阿佛洛狄忒}本人是一位行奸淫的女神。因为她有时以\Person{阿瑞斯}为情人，有时又以\Person{安喀塞斯}为情人，又有时以\Person{阿多尼斯}为情人，她还哀悼他的死亡，感到缺少爱人。他们还说她甚至下到\Place{冥府}去从\Person{珀耳塞福涅}那里买回\Person{阿多尼斯}。陛下，您可曾见过比这更愚蠢的事，竟将一个行奸淫、好哭泣哀号的女人奉为女神？\par

他们还描述\Person{阿多尼斯}是一位猎神，他惨遭横死，被野猪伤害，在困苦中无力自救。那么一个行奸淫、好打猎、又会死的人，又如何会关心人类呢？\par

所有这一切以及许多类似性质，甚至更为可耻和令人反感的事，陛下，希腊人都在他们的神祇身上讲述——这些事既不应说出来，也不应片刻记起。因此，人类从他们的神祇那里得到冲动，行一切不法、残暴和不敬之事，以他们可怕的行为玷污了大地和空气。\par

12. 埃及人又比这些民族更加愚昧无知，偏离正道比所有民族更远。因为他们不满足于加色丁人和希腊人的敬拜对象，此外还引入野兽作为神，包括陆地和水中动物、植物和草药；他们被一切疯狂和残暴所玷污，比地上所有民族更深。\par

因为他们起初敬拜\Person{伊西斯}，她将\Person{奥西里斯}作为兄弟和丈夫。\Person{奥西里斯}被自己的兄弟\Person{提丰}所杀；因此\Person{伊西斯}同她的儿子\Person{荷鲁斯}逃往\Place{叙利亚}的\Place{比布鲁斯}避难，为\Person{奥西里斯}痛苦哀哭，直到\Person{荷鲁斯}长大并杀死\Person{提丰}。因此\Person{伊西斯}既无力帮助自己的兄弟和丈夫；\Person{奥西里斯}被\Person{提丰}杀害时也不能自卫；\Person{提丰}这个杀兄弟者，在即将死于\Person{荷鲁斯}和\Person{伊西斯}之手时，也找不到自救之法。尽管他们因这些不幸暴露了真相，但单纯的埃及人却相信他们就是真神，他们甚至不满足于这些或其他民族的众神，还引入野兽作为神。因为有些人敬拜羊，有些人敬拜山羊；另一个部族敬拜公牛和猪；又有人敬拜乌鸦和鹰，以及秃鹫与雕；还有人敬拜鳄鱼；有些人敬拜猫和狗，以及狼和猿，以及龙和蝮蛇；另外有人敬拜洋葱、大蒜、荆棘和其他受造物。这些可怜的人却看不出这一切都是完全无力的。因为虽然他们看见自己的神被其他部族的人吃掉，被焚烧作为祭品，被宰杀作为牺牲，逐渐腐朽，他们仍不明白它们不是神。\par

13. 因此，埃及人、加色丁人和希腊人犯了极大的错误，竟将这样的一些存在奉为神明，为他们制作像，并将哑默无知的偶像神化。\par

我诧异他们看见自己的神被工匠锯开、劈砍、截断，又因年代久远而老化、破碎，且由金属铸成，却仍未看出它们不是神。\par

因为它们无力照管自己的安全，又如何能为人预作筹谋呢？\par

此外，加色丁人、希腊人和埃及人的诗人和哲学家，虽然想借他们的诗篇和著作来夸耀本国的神祇，反倒暴露出他们的羞耻，将之公诸世人。因为人的身体虽由许多部分组成，却不抛弃任何肢体，反而在所有肢体中保持完整的统一，自身和谐，那么在\OriginalTerm{天主}{God}的性体中，又怎能有如此大的分歧与不和谐呢？\par

因为如果在众神之间曾有本性的统一，那么一位神就不应追逐、杀害或伤害另一位。如果众神被其他神追逐、杀害、劫持、被雷击，那么就不再有任何本性的统一，而是各有分歧，尽是恶念。因此他们当中没有一位是神。那么显然，陛下，他们关于神性的一切论述都是错误的。\par

但是，希腊人的明智博学之士如何没有观察到，他们为自己立法，反被自己的律法所审判？如果律法是公义的，他们的神就全然不义，因为他们违法犯纪，互相残杀，施行巫术，行奸淫、偷窃和男色之事。如果他们做这些事是对的，那么律法就是不义的，因为它与神相反。然而事实上，律法是良善而公义的，称赞善事，禁止恶事。但他们神的行径却违反律法。因此他们的神是违法者，都当受死罪；而那些引人信奉这种神的人是不虔敬的。如果关于他们的故事是神话，那么这些神就不过是虚名；如果故事基于自然，那么行这些事并受这些苦的就不再是神；如果故事是寓言，那就只是神话，别无其他。\par

那么，陛下，已经表明所有这些多神的崇拜对象都是谬误和灭亡的作品。因为把可见而不可见的东西称为神是不对的；人应当敬畏那不可见、全见、全造的\OriginalTerm{天主}{God}。\par

14. 那么，大王，让我们也转向犹太人，看看他们对于天主的观点有何真理。因为他们是亚巴郎、依撒格和雅各伯的后裔，迁居到了埃及。天主以强有力的手和伸展的臂膀，藉着他们的立法者梅瑟将他们从那里领出；并用许多奇事和征兆向他们彰显了自己的大能。然而，他们却显出顽固和忘恩负义，屡次事奉异邦人的偶像，并杀害了派到他们那里的先知和义人。及至\OriginalTerm{天主子}{Son of God}乐意降临世间，他们竟以放肆的暴行接待他，将他出卖给罗马总督比拉多；毫不顾念他的善行和在他们中间所行的无数奇迹，竟要求处以十字架的死刑。\par

他们因自己的悖逆而灭亡；因为他们至今仍敬拜唯一的全能天主，但不是按照真知。他们否认\OriginalTerm{基督}{Christ}是\OriginalTerm{天主子}{Son of God}；他们与异教徒颇为相似，尽管他们似乎接近了他们所远离的真理。关于犹太人，就说到这里。\par

15. 基督徒的根源来自\OriginalTerm{主耶稣基督}{Lord Jesus Christ}。他被\OriginalTerm{圣神}{Holy Spirit}承认为至高天主之子，为救赎世人从天降下。由一位纯洁的童贞女所生，未经受孕，无玷无染，他取了肉身，在人间显示自己，为将世人从对多神的迷途召回。他完成了奇妙的\OriginalTerm{救恩计划}{dispensation}之后，自愿选择了十字架上的死亡，成就了庄严的救恩计划。三日后，他\OriginalTerm{复活}{resurrection}并\OriginalTerm{升天}{ascension}。大王，若您愿意阅读，可从他们所称的圣\OriginalTerm{福音}{gospel}书中判断他临在的光荣。他有十二位门徒，在他升天后，他们走遍世界各省，宣扬他的伟大。例如，其中一位遍历我们周围的国家，宣讲真理的教义。由此，凡仍遵守他们宣讲所吩咐的正义的人，被称为基督徒。\par

这些人比地上所有民族更找到了真理。因为他们通过\OriginalTerm{独生子}{the only-begotten Son}和\OriginalTerm{圣神}{Holy Spirit}认识\OriginalTerm{万物的创造者和塑造者}{Creator and Fashioner of all things}；除他外，他们不敬拜任何别的神。他们将主耶稣基督自己的诫命刻在心上，并遵守它们，期待着\OriginalTerm{死人的复活}{resurrection of the dead}和\OriginalTerm{来世的生命}{life in the world to come}。他们不犯奸淫，不犯邪淫，不作伪证，不贪图别人的东西；他们孝敬父母，爱邻人；他们公正审判，决不愿己所不欲加诸于人；他们劝诫伤害他们的人，并努力赢取他们为朋友；他们乐于对仇人行善；他们温柔易于被恳求；他们戒绝一切非法的言谈和一切不洁；他们不轻视寡妇，不压迫孤儿；有者无偿施舍以供养无者。\par

若看见陌生人，他们便接他到家中，像对待自己的弟兄一样为他欢喜；因为他们自称为弟兄，不是按肉身，而是按\OriginalTerm{心神}{spirit}。\par

他们随时准备为\OriginalTerm{基督}{Christ}舍命；因为他们坚定不移地遵守他的诫命，按照主天主所吩咐的，过圣洁而正义的生活。\par

他们每时每刻为主感恩，为一切饮食和其他恩赐。\par

16. 的确，这就是\OriginalTerm{真理之路}{way of the truth}，那些行在路上的人被引向通过\OriginalTerm{基督}{Christ}所应许的来世中的\OriginalTerm{永恒国度}{everlasting kingdom}。大王，为使您知道我说这些话并非出于私意，若您屈尊查阅基督徒的著作，您会发现我所说的没有超出真理。因此，您的儿子理解得正确，他受教事奉\OriginalTerm{永生天主}{living God}并要在那\OriginalTerm{将要来临的世代}{age that is destined to come upon us}中得救，这也是公义的。因为基督徒的言语和行为伟大而奇妙；他们说的不是人的话，而是天主的话。但其余的民族却迷途并欺骗自己；因为他们行走在黑暗中，像醉酒的人一样跌跌撞撞。\par

17. 大王，我对你说的这番话就到此为止，这是由\OriginalTerm{真理}{Truth}在我心中所指示的。因此，愿你们愚昧的智者们停止对主空洞的妄谈；因为敬拜天主创造者，并倾听他\OriginalTerm{不朽的圣言}{incorruptible words}，对你们是有益的，好使你们能逃脱定罪和惩罚，并成为\OriginalTerm{永生}{life everlasting}嗣业的继承者。\par

\SourceRecord{Translated by D.M. Kay.；Ante-Nicene Fathers；Vol. 9.；Edited by Allan Menzies.；Buffalo, NY: Christian Literature Publishing Co.,；1896.；<http://www.newadvent.org/fathers/1004.htm>.}
